\documentclass{article}

% Language setting
% Replace `english' with e.g. `spanish' to change the document language
\usepackage[english]{babel}

% Set page size and margins
% Replace `letterpaper' with `a4paper' for UK/EU standard size
\usepackage[letterpaper,top=2cm,bottom=2cm,left=3cm,right=3cm,marginparwidth=1.75cm]{geometry}

% Useful packages
\usepackage{amsmath}
\usepackage{graphicx}
\usepackage[colorlinks=true, allcolors=blue]{hyperref}
\usepackage[T2A]{fontenc}
\usepackage[utf8]{inputenc}
\usepackage[russian,english]{babel}
\usepackage{graphicx}

\begin{document}

\begin{center}
{\Huge Примерные задачи с физической интерпретацией\\}
\end{center}


{\Large \textbf{1) Простой перенос примеси}}

\[
\scalebox{1.3}{$u_t + a\,u_x = 0$}
\]

\noindent\textbf{Физическая интерпретация:}  
Дымо-газовая струя переносится постоянным ветром с постоянной скоростью.\\


{\Large \textbf{2) Адвекция-диффузия с постоянным коэффицентом}}

\[
\scalebox{1.3}{$u_t + a\,u_x = vu_x_x$}
\]

\noindent\textbf{Физическая интерпретация:}  
Выброс вещества в реке, течение уности загрязнение вниз, а турбулентность размывает его\\

{\Large \textbf{3) Адвекция-диффузия с переменным коэффицентом}}

\[
\scalebox{1.3}{$u_t + (a(x)\,u_x)_x = (v(x)u_x)_x$}
\]

\noindent\textbf{Физическая интерпретация:}  
Ветер усиливается на холме и ослабевает в низине, и примесь переносится неравномерно\\

{\Large \textbf{4) Движущийся источник\textsuperscript{*}}}


\[
\scalebox{1.3}{$u_t + a\,u_x = v\,u_x_x + S(x-c\,t)$}
\]

\noindent\textbf{Физическая интерпретация:}  
Автобус едет по дороге и выделяет выхлоп: концентрация газа формирует шлеф, который движется вместе с автобусом\\

{\Large \textbf{5) Переодическая среда\textsuperscript{*}}}

\[
\scalebox{1.3}{$\ensuremath{u_t + a(x)\,u_x = 0,\ \text{при } a(x + L) = a(x)}$}
\]
так же имеет место консервативная форма:
\[
\scalebox{1.3}{$\ensuremath{u_t + (a(x)\,u)_x = 0,\ \text{при } a(x + L) = a(x)}$}
\]

\noindent\textbf{Физическая интерпретация:}  
Воздух движется над местностью, где через какое-то количество расстояния чередуется лес и открытая местносность\\

{\Large \textbf{6) Сложные начальные условия}}

\[
\scalebox{1.3}{$u(x,0) = \sum A_i \cos(k_i x) + e^{-x^2}$}
\]


\noindent\textbf{Физическая интерпретация:}  
Атмосферный шум содержит набор частот, численная схема должна корректно передавать спектр при переносе\\


{\Large \textbf{7) Разрывная задача}}

\[
\scalebox{1.3}{$
u(x,0) =
\begin{cases}
1, & x < 0,\\[2mm]
0, & x \ge 0.
\end{cases}
$}
\]


\noindent\textbf{Физическая интерпретация:}  
Резкий выброс газа из балончика, концентрация скачом переходит от высокой к нулевой\\


{\Large \textbf{8) Адвекция с линейным источником\textsuperscript{*}\textsuperscript{*}}}

\[
\scalebox{1.3}{$u_t + a\, u_x = \sigma(x,t)\, u$}
\]

\noindent\textbf{Физическая интерпретация:}  
Рост (\(\sigma>0\)) или затухание (\(\sigma<0\)) переносимого вещества.  
Химическая реакция, радиоактивный распад, рост биомассы\\


{\Large \textbf{9) Адвекция с адаптивной сеткой}}

\[
\scalebox{1.3}{$u_t + a\, u_x = 0$}
\]

\noindent\textbf{Физическая интерпретация:}  
Сетка плотная там, где решение меняется резко (фронт, скачок), и редкая там, где профиль гладкий.  
Пример: поток загрязнения в реке с быстрым изгибом и спокойной областью.\\


{\Large \textbf{10) Слияние нескольких волн / фронтов}}

\[
\scalebox{1.3}{$u_t + a\, u_x = 0, \quad u(x,0) = u_1(x) + u_2(x)$}
\]

\noindent\textbf{Физическая интерпретация:}  
Два фронта (например, два загрязнения или два тепловых возмущения) распространяются в среде и могут встретиться. В месте встречи их амплитуды суммируются\\


{\Large \textbf{11) Перенос на круге / замкнутой области\textsuperscript{*}}}

\[
\scalebox{1.3}{$u_t + a\, u_x = 0, \quad x \in [0,L], \quad u(0,t) = u(L,t)$}
\]

\noindent\textbf{Физическая интерпретация:}  
Циркуляция атмосферы, кольцевой канал или поток по трубе кольцевого сечения. Волна движется по замкнутой области и возвращается к исходной точке, создавая повторяющееся поведение.\\


{\Large \textbf{12) Моделирование "текучей границы"}}

\[
\scalebox{1.3}{$u_t + a(u,x)\, u_x = 0$}
\]

\noindent\textbf{Физическая интерпретация:}  
Скорость переноса зависит от локального значения поля \(u\) и координаты \(x\). Примеры: река, где скорость потока зависит от концентрации загрязнения; транспорт биомассы; слабонелинейные потоки жидкости или газа. Фронт может ускоряться или замедляться в зависимости от локальной плотности.\\


\end{document}
